
Input to the Recharge (RCHA) Package is read from the file that has type ``RCH6'' in the Name File.  Any number of RCH6 packages can be specified for a single groundwater flow model.  The package type in the GWF Model name file PACKAGES block is always ``RCH6'' regardless of whether list-based or layer array-based input is used; the READASARRAYS option within the package input file activates layer array-based input.

Layer array-based input for recharge provides a similar approach for providing recharge rates as previous MODFLOW versions.  Layer array-based input for recharge can be used only with the DIS and DISV Packages.  Layer array-based input for recharge cannot be used with the DISU Package.

\mf will allow recharge in cells that also include a lake connection.  Refer to the ``\hyperref[sec:lakrchet]{Use of the Lake Package with RCH and EVT Packages}'' section for additional details when the RCHA and LAK packages are active in the same cell.

When layer array-based input is used for recharge, the DIMENSIONS block should not be specified.  The array size is determined from the model grid. 

\vspace{5mm}
\subsubsection{Structure of Blocks}
\vspace{5mm}

\noindent \textit{FOR EACH SIMULATION}
\lstinputlisting[style=blockdefinition]{./mf6ivar/tex/gwf-rcha-options.dat}
\vspace{5mm}
\noindent \textit{FOR ANY STRESS PERIOD}
\lstinputlisting[style=blockdefinition]{./mf6ivar/tex/gwf-rcha-period.dat}
\packageperioddescriptionarray{recharge}

\vspace{5mm}
\subsubsection{Explanation of Variables}
\begin{description}
\input{./mf6ivar/tex/gwf-rcha-desc.tex}
\end{description}

\vspace{5mm}
\subsubsection{Example Input File}
\lstinputlisting[style=inputfile]{./mf6ivar/examples/gwf-rcha-example.dat}

